\documentclass{article}
\usepackage[UTF8]{ctex}
\usepackage{amsmath,amssymb,amsfonts}
\usepackage{enumitem}
\usepackage{geometry}
\geometry{a4paper, margin=1in}
% === 错题本排版优化（不参与题目标识） ===
\setlength{\parindent}{0pt}
\setlength{\parskip}{0.35em plus 0.10em minus 0.08em}
\setlength{\emergencystretch}{3em}
\setlist{itemsep=0.25em, topsep=0.40em, parsep=0pt}
\allowdisplaybreaks[2]
% === 错题本排版优化结束 ===
\begin{document}
\section*{中国科学院大学2025年数学分析真题}
\section*{17．中国科学院大学2025年数学分析试卷}

\subsection{题目1}
1．计算下列极限。

\begin{itemize}
  \item[（1）] \(\lim _{x \rightarrow 1} x^{\frac{1}{1-x}}\) ．
  \item[（2）] \(\lim _{n \rightarrow \infty} \frac{1}{n^{4}}\left(1+2^{3}+\cdots+n^{3}\right)\) ．
\end{itemize}

\subsection{题目2}
2．用闭区间套定理证明［0，1］不可数．

\subsection{题目3}
3．设 \(f(x)\) 在 \([0,1]\) 上连续可微，\(f(0)=0\) ．证明：

\[
\int_{0}^{1}\left|f^{\prime}(x) f(x)\right| \mathrm{d} x \leq \frac{1}{2} \int_{0}^{1}\left|f^{\prime}(x)\right|^{2} \mathrm{~d} x .
\]

\subsection{题目4}
4．求重积分

\[
\iint_{D}(x+y) \sin (x-y) \mathrm{d} x \mathrm{~d} y
\]

其中 \(D=\{(x, y) \mid 0 \leq x+y \leq \pi, 0 \leq x-y \leq \pi\}\) ．

\subsection{题目5}
5．求曲面积分

\[
\iint_{\Gamma}\left(y^{2}-x\right) \mathrm{d} y \mathrm{~d} z+\left(z^{2}-y\right) \mathrm{d} z \mathrm{~d} x+\left(x^{2}-z\right) \mathrm{d} x \mathrm{~d} y
\]

其中 \(\Gamma=\left\{(x, y, z) \mid z=1-x^{2}-y^{2}, z \geq 0\right\}\) ，方向取外侧．

\subsection{题目6}
6．设 \(f(x)\) 在 \(\left(-\frac{\pi}{2}, 0\right)\) 上可积或绝对可积，怎么将 \(f(x)\) 延拓为到 \((-\pi, \pi)\) 上，满足其傅里叶级数为

\[
\sum_{n=1}^{\infty} b_{2 n-1} \sin (2 n-1) x
\]

\subsection{题目7}
7．设 \(f(x)\) 在 \([0, \pi]\) 上可积，且 \(\int_{0}^{\pi} f(x) \mathrm{d} x=\pi\) ．求系数 \(c_{1}, c_{2}, \cdots, c_{n}\) ，使得

\[
\int_{0}^{\pi}\left[f(x)-\sum_{k=1}^{n} c_{k} \cos k x\right]^{2} \mathrm{~d} x
\]

最小，并求 \(\lim _{n \rightarrow \infty} \sum_{k=1}^{n} c_{k} \cos k x=F(x)\) 的表达式。

\subsection{题目8}
8．设 \(f(x)\) 三次连续可微，且

\[
f^{\prime}(0)=1, f^{\prime \prime}(0)=0, f^{\prime \prime \prime}(0)=-1 .
\]

设 \(f\left(a_{n}\right)=a_{n+1}\) ，且 \(\lim _{n \rightarrow \infty} a_{n}=0\) ．求 \(\lim _{n \rightarrow \infty} n\left(a_{n}\right)^{2}\) ．

\subsection{题目9}
9．设 \(\left\{a_{n}\right\}\) 单调收敛到 \(0, \alpha \in(0, \pi)\) ，证明：\(\sum_{n=1}^{\infty} a_{n} \cos n x\) 在 \([\alpha, 2 \pi-\alpha]\) 上一致收敛。

\subsection{题目10}
10．求 \(\sum_{n=1}^{\infty}\left(1+\frac{1}{2}+\cdots+\frac{1}{n}\right) x^{n}\) 的收敛域与和函数．
\end{document}
